\documentclass{article}
\usepackage[utf8]{inputenc}
\usepackage[acronym]{glossaries}

\makeglossaries

\newglossaryentry{latex}
{
    name=latex,
    description={Is a markup language specially suited 
    for scientific documents}
}

\newglossaryentry{maths}
{
    name=mathematics,
    description={Mathematics is what mathematicians do}
}

\newacronym{gcd}{GCD}{Greatest Common Divisor}

\newacronym{lcm}{LCM}{Least Common Multiple}

\title{How to create a glossary}
\author{ }
\date{ }

\begin{document}
\maketitle
\acrlong{gcd}, which is abbreviated \acrshort{gcd}. This process 
is similar to that used for the \acrfull{lcm}.

The \Gls{latex} typesetting markup language is specially suitable 
for documents that include \gls{maths}. 

\printglossary[type=\acronymtype]
\printglossary

\end{document}